\documentclass{article}
\usepackage{amsmath, amssymb, amsthm}
\usepackage{hyperref}
\usepackage{geometry}
\geometry{margin=1in}

\title{Erd\H{o}s Problem \#1189}
\date{Last updated: 2026-04-04}
\author{Source: erdosproblems.com}

\begin{document}
\maketitle

\noindent\textbf{Status:} OPEN \quad \textbf{No prize}

\noindent\textbf{Tags:} number theory, covering systems

\medskip
\noindent\textbf{Problem Statement:}

Call a set of distinct integers $1<n_1<\cdots<n_k$ a covering set if there is a choice of $a_i\pmod{n_i}$ for $1\leq i\leq k$ such that every integer satisfies at least one of these congruences. A set is an irreducible covering set if no proper subset is a covering set. 

How many irreducible covering sets of size $k$ are there?

What is the minimum and maximum that $n_k$ can be?

Determine or estimate $\max \sum\frac{1}{n_i}$, where the maximum ranges over all irreducible covering sets of size $k$.

Are there infinitely many $n$ such that the divisors of $n$ (which are $>1$) form an irreducible covering set?

\bigskip
\noindent\textbf{Remarks:}

This is subtly different from the (more usually studied) notion of minimal covering system: an irreducible covering set corresponds to the set of moduli of a minimal covering system, but not necessarily vice versa.

The divisors of $12$ (which are $>1$) are an example of an irreducible covering set. Sun \cite{Su07} proved more generally that, for any odd prime $p$, the divisors (which are $>1$) of $2^{p-1}p$ form an irreducible covering set, settling the final question.

Simpson \cite{Si85} proved $n_k\leq 2^{k-1}$.

Let $I(k)$ count the number of irreducible covering sets of size $k$. Balister, Bollob\'{a}s, Morris, Sahasrabudhe, and Tiba \cite{BBMST24} proved that\[I(k) \leq \exp\left((c+o(1))\frac{k^{3/2}}{(\log k)^{1/2}}\right)\]by proving that the right-hand side is an asymptotic for the number of minimal covering systems on $k$ moduli.

It is trivial that $\sum\frac{1}{n_i}>1$ for any irreducible covering set.

See also [1188].

\bigskip
\noindent\textbf{References:}

[BBMST24] Balister, Paul and Bollob\'{a}s, B\'ela and Morris, Robert and
Sahasrabudhe, Julian and Tiba, Marius, The structure and number of {E}rd\H os covering systems. J. Eur. Math. Soc. (JEMS) (2024), 75--109.
 
 [Si85] Simpson, R. J., Regular coverings of the integers by arithmetic progressions. Acta Arith. (1985), 145--152.
 
 [Su07] Sun, Zhi-Wei, On covering numbers. (2007), 443--453.

\bigskip
\noindent\small{Source: \url{https://www.erdosproblems.com/1189}}
\end{document}
